% source: f4fd6529c0aa1a86b4cc1b3094c3893b27c87452
\chapter{Gram--Schmidt 正交化与 LLL 算法}
\label{ch:lll}

\section{本章为何重要}
本章常被称为格基约化的核心。原因很简单：一旦我们理解了格是什么，以及为什么某些基比另一些更好，就需要一种方法来度量并改进这种质量。Gram--Schmidt 正交化恰好提供了这一工具。它是 LLL、BKZ、Babai 最近平面算法~\cite{babai1986}以及许多其他格技术的概念基础。

\section{学习目标}
读完本章后，你应当能够：

\begin{itemize}
  \item 解释为什么有些格基是好的，而另一些是坏的；
  \item 理解正交性的含义；
  \item 从概念上描述 Gram--Schmidt 正交化做了什么；
  \item 理解 LLL 为什么反复使用 Gram--Schmidt 信息；
  \item 将基的质量与 LLL、BKZ 等后续算法联系起来。
\end{itemize}

\section{什么样的基是好基？}
我们从直观认识而非公式入手。考虑平面上的两组基。一组中，向量短且近乎垂直；另一组中，向量长且近乎平行。二者可能生成同一个格，但第二组处理起来要困难得多。

好基（good basis）是让格易于“导航”的基。坏基则把短向量隐藏起来，使格看上去比实际复杂得多。

\section{为什么近乎平行的向量是坏的}
假设我们有一组坏基，由两个几乎共线的向量组成。为了表示像 $(1,0)$ 这样简单的向量，可能需要用到基向量的非常大的整数组合。这意味着坐标表示变得不稳定且数值很大。在密码学中，这正是我们要避免的。

这就是格基约化旨在寻找更好的基的原因。

\section{正交是什么意思？}
两个向量 $u$ 与 $v$ 称为正交的，如果

\[
u \cdot v = 0.
\]

例如，$(1,0)$ 与 $(0,1)$ 是正交的；相反，$(2,1)$ 与 $(3,2)$ 不是。

\subsection{Sage 实验：点积}

\begin{lstlisting}[language=Python]
u = vector(QQ, [1, 0])
v = vector(QQ, [0, 1])
u.dot_product(v)
\end{lstlisting}

结果为 $0$，表明这两个向量正交。

\section{Gram--Schmidt 做了什么}
Gram--Schmidt 过程不改变格。从根本意义上说，它甚至不改变基。它所做的是分析一组基，把每个向量拆分为两部分：

\begin{itemize}
  \item 沿先前各方向的分量；
  \item 指向真正新方向的新分量。
\end{itemize}

通过这种方式，该过程分离出每个向量中贡献新几何信息的那一部分。所得向量彼此正交，并揭示了基中蕴含的真实方向。

这就是为什么 Gram--Schmidt 常被描述为一种剥离旧方向、只保留新方向的方法。

\section{投影公式}
若 $u$ 与 $v$ 是向量，则 $v$ 在 $u$ 上的投影为

\[
\operatorname{proj}_u(v)=\frac{u\cdot v}{u\cdot u}u.
\]

Gram--Schmidt 过程从 $v$ 中减去这一投影，留下 $v$ 中与 $u$ 正交的部分。

\section{Gram--Schmidt 构造}
第一个向量保持不变：

\[
b_1^* = b_1.
\]

随后的每个向量都通过减去其在先前各方向上的投影来调整：

\[
b_i^* = b_i - \sum_{j<i} \mu_{i,j} b_j^*,
\]

其中系数 $\mu_{i,j}$ 的选取使得新向量与先前各向量正交。

\begin{figure}[H]
\centering
\begin{tikzpicture}[scale=1.15]
  \draw[pqcgray!30, thin] (-1,0) -- (3.5,0);
  \draw[pqcgray!30, thin] (0,-0.5) -- (0,2.5);
  \draw[basisvec] (0,0) -- (3,1);
  \node[pqcblue, font=\small] at (3.2,0.85) {$b_1$};
  \draw[altvec] (0,0) -- (1,2);
  \node[pqcteal, font=\small] at (0.75,2.2) {$b_2$};
  \node[circle, fill=pqcgray, inner sep=0pt, minimum size=2.5pt] at (1.5,0.5) {};
  \node[pqcgray, font=\scriptsize] at (2.15,0.28) {$\mathrm{proj}_{b_1}(b_2)$};
  \draw[shortvec] (1.5,0.5) -- (1,2);
  \node[pqcorange, font=\small] at (0.5,1.4) {$b_2^{*}$};
\end{tikzpicture}
\caption{$b_2$ 相对于 $b_1$ 的 Gram--Schmidt 分解，其中 $b_1=(3,1)$，$b_2=(1,2)$。圆点标出 $\mathrm{proj}_{b_1}(b_2)=(1.5,0.5)$，即 $b_2$ 沿 $b_1$ 方向的分量。橙色向量 $b_2^{*}=b_2-\mathrm{proj}_{b_1}(b_2)$ 与 $b_1$ 正交：它是 $b_2$ 在 $b_1$ 之外贡献的真正新的几何信息。}
\label{fig:gram-schmidt}
\end{figure}

这一公式可以直接推广为一个算法：逐个处理基向量，对每个向量减去它在所有先前已正交化方向上的投影。

\begin{algorithm}[H]
\caption{Gram--Schmidt 正交化}
\label{alg:gram-schmidt}
\begin{algorithmic}[1]
\Require 基 $b_1,\dots,b_n$
\Ensure 正交向量 $b_1^*,\dots,b_n^*$ 及系数 $\mu_{i,j}$
\State $b_1^* \gets b_1$
\For{$i = 2,\dots,n$}
    \State $b_i^* \gets b_i$
    \For{$j = 1,\dots,i-1$}
        \State $\mu_{i,j} \gets \dfrac{b_i \cdot b_j^*}{b_j^* \cdot b_j^*}$
        \State $b_i^* \gets b_i^* - \mu_{i,j}\, b_j^*$
    \EndFor
\EndFor
\State \Return $(b_1^*,\dots,b_n^*)$ 及 $(\mu_{i,j})$
\end{algorithmic}
\end{algorithm}

\section{Sage 实验：Gram--Schmidt}

\begin{lstlisting}[language=Python]
B = matrix(QQ, [[2, 1], [1, 3]])
B.gram_schmidt()
\end{lstlisting}

输出是由原始基导出的一组正交向量。它们本身不一定是格向量，但对分析极为有用。

\subsection{验证}

\begin{lstlisting}[language=Python]
u = vector(QQ, [2, 1])
v = vector(QQ, [-1, 2])
u.dot_product(v)
\end{lstlisting}

点积应为 $0$，表明这两个向量正交。

\section{为什么 Gram--Schmidt 不改变格}
这是一个微妙但重要的问题。Gram--Schmidt 向量通常不是整数向量，甚至可能根本不在格中。它们是分析工具，而不是新的格生成元。它们的作用是揭示格基的几何结构，而不是取代格基。

\section{Gram--Schmidt 长度}
Gram--Schmidt 向量的长度

\[
\|b_i^*\|,
\]

是衡量基质量的关键指标。如果这些长度平缓递减并保持相当均衡，则该基是好的；如果某个向量变得极短而另一个变得极长，则该基质量较差。

\section{正交性缺陷}
正交性缺陷（orthogonality defect）是一个有用的质量度量。其一种形式为

\[
\mathrm{OD}(B)=\frac{\prod_i \|b_i\|}{\det(B)}.
\]

若基向量完全正交，则缺陷为 $1$。缺陷越大，基越差。

\section{Hadamard 比}
另一个标准量是 Hadamard 比：

\[
\mathrm{HR}(B)=\left(\frac{\det(B)}{\prod_i \|b_i\|}\right)^{1/n}.
\]

它的取值介于 $0$ 与 $1$ 之间，越接近 $1$ 表示基的质量越好。

\section{LLL 为什么反复使用 Gram--Schmidt}
LLL 算法~\cite{lll1982}建立在这样一个思想之上：通过反复检查 Gram--Schmidt 信息，可以逐步改进一组基。如果当前基约化得不够充分，LLL 就交换向量、重新缩放，并重新计算正交化方向，直到基的质量可以接受为止。

因此，Gram--Schmidt 并非一个旁枝概念：它是最重要的格基约化算法之一背后的引擎。该算法在两种操作之间交替进行：\emph{尺寸约化}（size reduction），利用系数 $\mu_{i,j}$ 使每个 $b_i$ 尽可能接近先前各向量张成的空间，同时不改变格；以及 \emph{Lov\'asz 条件}，这是对相邻 Gram--Schmidt 长度的一项检验，用以决定两个相邻向量是否应当交换。

\begin{algorithm}[H]
\caption{LLL 格基约化}
\label{alg:lll}
\begin{algorithmic}[1]
\Require 基 $b_1,\dots,b_n$；约化参数 $\delta \in (1/4,1)$，通常取 $\delta = 3/4$
\Ensure 一组与 $b_1,\dots,b_n$ 生成同一个格的 LLL 约化基
\State 通过 Gram--Schmidt（算法~\ref{alg:gram-schmidt}）计算 $b_1^*,\dots,b_n^*$ 与 $\mu_{i,j}$
\State $k \gets 2$
\While{$k \le n$}
    \For{$j = k-1$ 递减至 $1$} \Comment{对 $b_k$ 关于每个先前向量做尺寸约化}
        \If{$|\mu_{k,j}| > 1/2$}
            \State $b_k \gets b_k - \mathrm{round}(\mu_{k,j})\, b_j$
            \State 重新计算 $b_k^*$ 与 $\mu_{k,\cdot}$
        \EndIf
    \EndFor
    \If{$\|b_k^*\|^2 \ge (\delta - \mu_{k,k-1}^2)\, \|b_{k-1}^*\|^2$} \Comment{Lov\'asz 条件成立}
        \State $k \gets k+1$
    \Else
        \State 交换 $b_k \leftrightarrow b_{k-1}$
        \State 重新计算 $b_1^*,\dots,b_n^*$ 与 $\mu_{i,j}$
        \State $k \gets \max(k-1,\,2)$
    \EndIf
\EndWhile
\State \Return $(b_1,\dots,b_n)$
\end{algorithmic}
\end{algorithm}

每次交换都会严格减小一个由 Gram--Schmidt 长度构造的势函数，因此循环必然终止；在实践中，即使在基于格的密码体制所用的数百维情形下，它也只需少数几轮即可收敛。输出并不是可能的最短基，但它可证明地足够短——每个 $\|b_i\|$ 都能用真正的最短向量加以界定——这使得 LLL 成为任何正式格攻击之前的主力预处理步骤。

\section{与 LWE 的联系}
本章也有助于解释为什么格问题是困难的。在 LWE 中，秘密和误差隐藏在一个高维结构之中。如果攻击者能够得到相关格的一组好基，许多相关任务都会变得容易得多。基于格的密码的安全性正是建立在这样的假设之上：这样的基不易获得。

\section{实验}

\subsection{实验 1：正交性}

\begin{lstlisting}[language=Python]
u = vector(QQ, [1, 0])
v = vector(QQ, [0, 1])
u.dot_product(v)
\end{lstlisting}

\subsection{实验 2：比较两对向量}

\begin{lstlisting}[language=Python]
u = vector(QQ, [100, 99])
v = vector(QQ, [101, 100])
u.dot_product(v)
\end{lstlisting}

将结果与前一个实验进行比较。

\subsection{实验 3：在 Sage 中做 Gram--Schmidt}

\begin{lstlisting}[language=Python]
B = random_matrix(QQ, 3)
B.gram_schmidt()
\end{lstlisting}

\subsection{实验 4：正交性缺陷}
生成若干组随机基，并比较它们的正交性缺陷值。

\section{小结}
本章介绍了基质量的几何概念。Gram--Schmidt 正交化并不是一种新的格构造；它是一种分析基的方法，把旧方向与真正的新方向区分开来。好基短且近乎正交，而坏基长且近乎平行。这些思想是 LLL 及其他约化算法的直接前导。

下一章\emph{最短向量问题（SVP）}将把这些几何思想与一个具体的困难问题联系起来。此后，本书将转向噪声分布和带误差学习问题，解释为什么格问题与现代后量子密码的安全性如此紧密相关。

\printbibliography[heading=chapterrefs]
